\documentclass[11pt]{ctexart}
\usepackage{amsmath,amssymb}
\begin{document}
\section*{P204：完整当前整数时槽双层模型}
本版恢复历史所选02的整数时槽形式：保留式8同一$q$每时刻最多一船，
式11--14按完整小时计算处理量，式32--34每个二元岸电变量计一整小时。

两个明示定义修复为：原式5的泊位指标改为对每个$b$检查所选泊位长度；
原式32--34时间端点按小时起点$s_i\le t<d_i$解释，避免离泊端点重复收费。
按原论文PDF第10页图2将靠泊燃料区间修正为$d_i-s_i-T_i^E$，
修正第13页式26与图示的内部冲突；不重复计入锚泊燃料。
箱量取最近的正$\eta$倍数只修复数据相容性，逐船前后值另存；
其余ETA、截止时间、罚金、船舶、泊位、有限岸桥和能源数据不变。

\section*{集合、参数及变量}
$V,B,R$为船舶、泊位、船公司；$\mathcal T=\{0,\ldots,H-1\}$为小时起点，
$\mathcal Q=\{1,\ldots,Q\}$为岸桥数量选项，$H_i\subseteq\mathcal Q$为船$i$允许的数量。
$c_i,l_i,\eta,L_b,D_i,\sigma_i,ETA_i$为箱量、船长、每桥每小时处理率、
泊位长度、截止时间、延误罚金及参考到港时刻。
$e_i,\alpha_i,\beta_b\in\{0,1\}$为船舶岸电、脱硫及泊位岸电标志。
上层变量为$b_i\in B$、$0\le s_i\le d_i\le H$，
$x_{itq},y_{itb}\in\{0,1\}$。
$x_{itq}=1$表示小时$t$为船$i$配置$q$台岸桥；
$y_{itb}=1$为原式的泊位及排序时间标记。
下层变量为合法燃料航段速度$v_{ig}$、到港时间$a_i\ge0$、
$\gamma_{it}\in\{0,1\}$及由下式定义的非负时间、能耗和费用。
船舶次序、每船岸桥数量$q$和并行作业安排均由模型选择，$s_i,d_i$采用连续时间域。

\section*{上层目标及原离散可行域}
\[
\min(U_1,U_2,U_3),\quad
U_1=\sum_{i:e_i=1}\sigma_i(d_i-D_i)^+,\quad
U_2=\sum_{i:\alpha_i=1}\sigma_i(d_i-D_i)^+,\quad
U_3=\sum_i\sigma_i(d_i-D_i)^+.
\]
采用正权标量$2U_1+U_2+U_3$计算方案的综合目标值。
\begin{align}
a_i&\le s_i &&\forall i \tag{4}\\
l_i\sum_{t\in\mathcal T}y_{itb}&\le L_b &&\forall i,b \tag{5}\\
\sum_{t,b}y_{itb}&=1 &&\forall i \tag{6}\\
\sum_i y_{itb}&\le1 &&\forall t,b \tag{7}\\
\sum_i x_{itq}&\le1 &&\forall t,q \tag{8}\\
\sum_{t,b}b\,y_{itb}&=b_i &&\forall i \tag{9}\\
\sum_i\sum_{q\in H_i}q x_{itq}&\le Q &&\forall t \tag{10}\\
s_i+\sum_t\sum_{q\in H_i}x_{itq}&\le d_i &&\forall i \tag{11}\\
\eta\sum_t\sum_{q\in H_i}q x_{itq}&=c_i &&\forall i \tag{12}\\
(t+1)\sum_{q\in H_i}x_{itq}&\le d_i &&\forall i,t \tag{13}\\
t\sum_{q\in H_i}x_{itq}+H(1-\sum_{q\in H_i}x_{itq})&\ge s_i &&\forall i,t \tag{14}\\
d_i y_{itb}&\le s_j y_{jt'b}+H(1-y_{jt'b})
&&\forall i\ne j,b,t<t' \tag{15}\\
x_{itq},y_{itb}&\in\{0,1\}. && \tag{16}
\end{align}
没有另外把每船每时刻只能选一个$q$加入完整模型；候选算法自愿使用该可行子集。
$q\notin H_i$的无用$x$可置零，不减少任何目标有效选择。

\section*{下层全部目标、燃料及整小时岸电约束}
$G_i$列出船舶全部正距离合法燃料航段；未装脱硫器采用对应区域合法非HFO燃料，
装脱硫器时允许HFO。参数为$w_{ig}>0,P_{ig}>0,F^{main}_{ig}>0,
F^{aux}_{ig}\ge0,u_i,v^0_{ig},\underline v_{ig},\bar v_{ig}$。
港内燃料每小时成本$p_i^a$由相应燃料价格和辅机消耗率相乘给出。
岸电功率$\theta_i$、每整小时电价$P_t^E$、基础港费$\Psi_i^0$、
折扣比例$\lambda$和门槛$T^*$为输入。
固定上层后，各公司$r$最小化$\sum_{i:r(i)=r}C_i$：
\begin{align*}
a_i&=ETA_i+\sum_{g\in G_i}w_{ig}(v_{ig}^{-1}-(v^0_{ig})^{-1}),&
0&\le a_i\le s_i,\\
\underline v_{ig}&\le v_{ig}\le\bar v_{ig},&
\gamma_{it}&\in\{0,1\},\\
T_{ik}^E&=\sum_{t\in K_k:\,s_i\le t<d_i}
             \gamma_{it}\min(e_i\beta_{b_i},1),&
k&\in\{p,mp,op\},\\
T_i^E&=\sum_kT_{ik}^E,&
0&\le T_i^E\le d_i-s_i,\\
E_i^E&=\sum_k P_k^E\theta_i T_{ik}^E,&
C_i^{sail}&=\sum_gP_{ig}\big[F^{main}_{ig}(v_{ig}/u_i)^3+F^{aux}_{ig}\big]w_{ig}/v_{ig},\\
C_i^{anc}&=p_i^a(s_i-a_i),&
C_i^{berth}&=p_i^a(d_i-s_i-T_i^E),\\
\Psi_i&=\begin{cases}\Psi_i^0&T_i^E<T^*,\\
\lambda\Psi_i^0&T_i^E\ge T^*,\end{cases}&
C_i&=C_i^{sail}+C_i^{anc}+C_i^{berth}+E_i^E+\Psi_i.
\end{align*}
$K_k$为输入给定的峰、平、谷时段集合；逐小时电价向量与上述分组相同。
每个有效$\gamma$贡献一小时；不存在按交叠长度收费的连续边缘小时。
无效时段/不具备岸电条件时的$\gamma$可以置零，不影响目标或可行域。
公司之间没有新增共享岸电功率约束。双层要求下层是完整反应问题的全局最优解。

\section*{求解和独立证书}
上层使用多种船舶次序的有限资源整数时槽构造，允许逐时变化岸桥数量；
候选的$x,y,\gamma$逐项保存，未列的二元变量均为零。
证书对每个$t,q,b$复算原式4--16、32--36、全部上下界及成本，
不以候选生成规则替代完整约束核验。

令$\tau_{ig}=w_{ig}/v_{ig}$，
$K_{ig}=P_{ig}F^{main}_{ig}w_{ig}^3/u_i^3>0$，
$D_{ig}=P_{ig}F^{aux}_{ig}-p_i^a$。
航速项等价为$\sum_g(K_{ig}/\tau_{ig}^2+D_{ig}\tau_{ig})$，
保留全部速度箱约束及$0\le ETA_i-\sum_gw_{ig}/v^0_{ig}+\sum_g\tau_{ig}\le s_i$。
严格正二阶导数$6K_{ig}/\tau_{ig}^4$使可行KKT证书足以证明全局航速最优。
对每个可行岸电小时数$k$，选择增量成本$P_t^E\theta_i-p_i^a$最低的$k$个整小时，
枚举全部$k$并计算港费阈值；这给出全二元岸电域的精确最优反应。
两部分可分离，按船求解得到固定上层下的公司全局最优反应。
求解入口不调用Gurobi，单调时钟上限3600秒，乘子采用有限解析括区间和最多160次二分。
输出保守记为启发式双层可行点，界/gap为空；若延迟向量恰为零，
另给非负目标达到零的证书，不以此宣称全部Pareto前沿或下层总成本跨上层全局最优。
\end{document}
